%% toc: P2 15 15.5,15.6,15.7
\section{Proofs for Horn-SAT, 2-SAT and propositional decision complexity}
\label{sec:annex-p02-w02-p02-019-sat-procedures}

Throughout, clause lists, \(A_S\), \(n\), \(m\), \(N\), literal codes and the
random-access machine are as in
Convention~\ref{def:p02-w02-p02-019-clause-input}.  Satisfaction of a clause
is used in the form of Lemma~\ref{lem:prop-clause-satisfaction}: \(v\)
satisfies \(C\) iff some literal of \(C\) is true under \(v\).

\begin{annexitem}{annex:p02-w02-p02-019-hc}{thm:p02-w02-p02-019-horn-closure}{Forward closure}
\begin{proof}
Write \(F=\operatorname{Fwd}(S)\), and let \(\mathcal K\) be the family of
subsets of \(A_S\) closed under \(S\).

\emph{Clause 1.}  \(A_S\in\mathcal K\): the head of every definite clause of
\(S\) occurs in \(S\), so it belongs to \(A_S\).  Let \(C\) be a definite
clause of \(S\) with \(B(C)\subseteq F\).  For every \(X\in\mathcal K\),
\(F\subseteq X\), so \(B(C)\subseteq X\) and, \(X\) being closed,
\(\operatorname{hd}(C)\in X\).  Hence \(\operatorname{hd}(C)\) lies in every
member of \(\mathcal K\), that is, in \(F\).  So \(F\) is closed, and by its
definition as an intersection it is contained in every member of
\(\mathcal K\).

\emph{Clause 2.}  Let \(v\) satisfy \(S\) and put
\(X_v=\{q\in A_S: v(q)=\mathbf{T}\}\).  Let \(C\) be a definite clause of \(S\)
with \(B(C)\subseteq X_v\).  Every negative literal \(\neg q\) of \(C\) has
\(q\in B(C)\), so \(v(q)=\mathbf{T}\) and \(\neg q\) is false under \(v\).  Since
\(v\) satisfies \(C\), some literal of \(C\) is true; the only literal of \(C\)
that is not negative is its head, so \(v(\operatorname{hd}(C))=\mathbf{T}\) and
\(\operatorname{hd}(C)\in X_v\).  Thus \(X_v\in\mathcal K\), and by clause~1,
\(F\subseteq X_v\): every atom of \(F\) is true under \(v\).

\emph{Clause 3.}  Let \(C\) be a definite clause of \(S\).  If
\(B(C)\subseteq F\), then \(\operatorname{hd}(C)\in F\) by clause~1, so the
head is true under \(v_F\) and \(v_F\) satisfies \(C\).  Otherwise some
\(q\in B(C)\) is not in \(F\); then \(v_F(q)=\mathbf{F}\), the literal \(\neg q\)
of \(C\) is true, and \(v_F\) satisfies \(C\).  Let \(C\) be a negative clause.
Its literals are exactly the \(\neg q\) with \(q\in B(C)\).  So \(v_F\)
satisfies \(C\) iff some \(q\in B(C)\) has \(v_F(q)=\mathbf{F}\), iff some
\(q\in B(C)\) is not in \(F\), iff \(B(C)\not\subseteq F\).  For the empty
clause, \(B(C)=\varnothing\subseteq F\) and the clause is false, as it must
be.

\emph{Clause 4.}  Suppose \(S\) contains a negative clause \(C\) with
\(B(C)\subseteq F\), and let \(v\) be any valuation satisfying \(S\).  By
clause~2 every \(q\in B(C)\) is true under \(v\), so every literal of \(C\) is
false and \(v\) does not satisfy \(C\), a contradiction.  So \(S\) is
unsatisfiable.  Conversely, if every negative clause \(C\) of \(S\) has
\(B(C)\not\subseteq F\), clause~3 shows that \(v_F\) satisfies every definite
and every negative clause of \(S\); a Horn clause is one or the other, so
\(v_F\) satisfies \(S\).  The two directions together give the contradiction
condition.  If \(S\) is satisfiable and \(v\) satisfies \(S\), the atoms true
under \(v_F\) are those of \(F\), and they are true under \(v\) by clause~2.

\emph{Clause 5.}  If \(S\) has no positive unit clause, every definite clause
of \(S\) has a nonempty body, which is not a subset of \(\varnothing\); so
\(\varnothing\in\mathcal K\) and \(F=\varnothing\).  A negative clause with
\(B(C)\subseteq\varnothing\) has empty body and is the empty clause.  If there
is no empty clause, clause~4 gives that \(v_\varnothing\), the valuation
making every variable false, satisfies \(S\).  For the empty list,
\(A_S=\varnothing\), \(\mathcal K=\{\varnothing\}\) and \(F=\varnothing\);
there is no clause to falsify, so every valuation satisfies it.
\end{proof}
\end{annexitem}

\begin{annexitem}{annex:p02-w02-p02-019-hp}{thm:p02-w02-p02-019-horn-procedure}{Forward chaining is correct and linear}
\begin{proof}
If \(S\) contains an empty clause, step~(0) finds the first one, say \(C_j\),
and halts with ``unsatisfiable'', which is correct by
Theorem~\ref{thm:p02-w02-p02-019-horn-closure}, clause~4.  It inspected
\(j\) clauses, and the \(j-1\) clauses before \(C_j\) are nonempty, so
\(j-1\leq N\).  From now on assume that no clause is empty.

\emph{Step (1).}  Each \(C_j\) is a list without repetition, so after
step~(1): \(\mathit{cnt}[j]=|B(C_j)|\); \(\mathit{hd}[j]\) is the index of the
head if \(C_j\) is definite and \(0\) if it is negative (a Horn clause has at
most one positive literal, so \(\mathit{hd}[j]\) is written at most once); and
\(\mathit{occ}[q]\) lists, once each, the \(j\) with \(p_q\in B(C_j)\).

\emph{Invariants.}  At any moment after step~(1), let \(M\) be the set of
marked atoms and \(D\) the set of atoms already removed from \(Q\).  Before
step~(2) and after each complete iteration of step~(2) or of the loop in
step~(3), the following hold.
\begin{enumerate}
  \item[(a)] An atom is appended to \(Q\) only at the moment it is marked,
        and marks are never removed.  Hence each atom is appended at most
        once, \(D\subseteq M\), and \(Q\) contains exactly the atoms of
        \(M\setminus D\).
  \item[(b)] \(\mathit{cnt}[j]=|B(C_j)\setminus D|\) for every \(j\).
  \item[(c)] For every definite \(C_j\) with \(\mathit{cnt}[j]=0\),
        \(\operatorname{hd}(C_j)\in M\).
  \item[(d)] \(M\subseteq\operatorname{Fwd}(S)\).
\end{enumerate}
Before step~(2), \(M=D=\varnothing\), \(Q\) is empty and (b) is the
statement about step~(1).  For (c) at that moment, note that a clause with
\(\mathit{cnt}[j]=0\) then has empty body; it is a positive unit clause, and
step~(2) marks its head.  Step~(2) does not remove atoms, so (b) persists
through it, and after it (c) holds.  In step~(3), removing \(q\) from \(Q\)
adds \(q\) to \(D\); by (a), \(q\notin D\) before.  The loop then decrements
\(\mathit{cnt}[j]\) exactly once for each \(j\) with \(p_q\in B(C_j)\), so
(b) is preserved.  A counter changes only by these decrements; whenever one
becomes \(0\) for a definite clause, the head is marked at once if it was not
already marked.  Together with (c) before the iteration, this gives (c) after
it.  For (d), consider the moment an atom \(r\) is marked.  It is the head of
a definite clause \(C_j\) with \(\mathit{cnt}[j]=0\), so by (b)
\(B(C_j)\subseteq D\subseteq M\) at that moment.  By (d) before the mark,
\(B(C_j)\subseteq\operatorname{Fwd}(S)\).  Since \(\operatorname{Fwd}(S)\) is
closed (Theorem~\ref{thm:p02-w02-p02-019-horn-closure}, clause~1),
\(r\in\operatorname{Fwd}(S)\).

\emph{Termination and the closure.}  By (a), at most \(n\) atoms are ever
appended, so the loop of step~(3) runs at most \(n\) times, and each run
handles the finite list \(\mathit{occ}[q]\).  When the loop ends, \(Q\) is
empty, so \(D=M\) by (a).  Let \(C_j\) be definite with
\(B(C_j)\subseteq M\).  By (b), \(\mathit{cnt}[j]=|B(C_j)\setminus M|=0\), and
by (c) \(\operatorname{hd}(C_j)\in M\).  So \(M\) is closed under \(S\).  By
Theorem~\ref{thm:p02-w02-p02-019-horn-closure}, clause~1,
\(\operatorname{Fwd}(S)\subseteq M\), and with (d),
\(M=\operatorname{Fwd}(S)\).

\emph{Output.}  In step~(4), for a negative clause \(C_j\), (b) and \(D=M\)
give \(\mathit{cnt}[j]=0\) iff \(B(C_j)\subseteq\operatorname{Fwd}(S)\).  So
step~(4) outputs ``unsatisfiable'' iff some negative clause has its body
inside \(\operatorname{Fwd}(S)\), iff \(S\) is unsatisfiable, by
Theorem~\ref{thm:p02-w02-p02-019-horn-closure}, clause~4.  Otherwise the
returned valuation makes true exactly the atoms of \(M=\operatorname{Fwd}(S)\);
it is \(v_{\operatorname{Fwd}(S)}\), which satisfies \(S\) by the same clause.

\emph{Cost.}  Step~(0) costs \(O(m)\).  Step~(1) creates arrays of total
length \(2n+2m\) in \(O(n+m)\) steps and reads every literal occurrence once,
doing \(O(1)\) work for it: \(O(N)\).  Step~(2) inspects \(m\) counters:
\(O(m)\).  In step~(3) each atom is removed at most once by (a); removing
\(q\) costs \(O(1)\) plus \(O(1)\) for each entry of \(\mathit{occ}[q]\).  The
entries of all the lists \(\mathit{occ}[q]\) correspond one to one to the
negative literal occurrences of \(S\), so step~(3) costs
\(O(n)+O(N)\), and each negative occurrence is the source of at most one
decrement.  Step~(4) costs \(O(m)\) for the goal check and \(O(n)\) to write
the valuation.  The total is \(O(m+n+N)\).  With no empty clause,
\(m\leq N\) and \(n\leq N\) by
Convention~\ref{def:p02-w02-p02-019-clause-input}, so the total is \(O(N)\).
\end{proof}
\end{annexitem}

\begin{annexitem}{annex:p02-w02-p02-019-scc}{thm:p02-w02-p02-019-scc}{Two-pass strong components}
\begin{proof}
For a call \(\mathrm{Visit}(u)\) of the first pass, its \emph{interval} is
the time from the moment \(u\) turns grey to the moment \(u\) turns black.

\emph{Fact 1 (nesting).}  Every vertex turns grey exactly once: a call is made
only on a white vertex, which it turns grey, and no vertex turns white again;
the outer loop calls \(\mathrm{Visit}\) on every vertex still white.  If
\(\mathrm{Visit}(w)\) is called during the interval of \(\mathrm{Visit}(u)\),
with \(w\neq u\), the call is made from a call that is itself active, and
calls return in the reverse order in which they were made.  So the interval of
\(w\) lies inside that of \(u\), and \(f(w)<f(u)\).

\emph{Fact 2 (only reachable vertices).}  Every vertex turned grey during the
interval of \(\mathrm{Visit}(u)\) is reachable from \(u\).  Indeed it is
turned grey by a call made through an edge from a vertex that turned grey
earlier in the same interval, or from \(u\) itself; induction on the order in
which vertices turn grey gives a path from \(u\).

\emph{Fact 3 (white paths).}  Suppose that, when \(u\) turns grey, there is a
path \(u=x_0\to x_1\to\cdots\to x_k=w\) whose vertices \(x_1,\ldots,x_k\) are
all white.  Then \(w\) turns grey during the interval of \(\mathrm{Visit}(u)\),
and if \(w\neq u\) then \(f(w)<f(u)\).  Suppose not, and let \(i\geq 1\) be
least such that \(x_i\) does not turn grey during the interval.  Then
\(x_{i-1}\) is \(u\) or turned grey during the interval, and by Fact~1 its
call runs inside the interval and examines the edge \(x_{i-1}\to x_i\) before
returning.  The vertex \(x_i\) was white when \(u\) turned grey and does not
turn grey during the interval, so it is white at that examination, and
\(\mathrm{Visit}(x_i)\) is called inside the interval.  This contradicts the
choice of \(i\).  The inequality follows by Fact~1.

\emph{Fact 4.}  A path between two vertices of one strong component \(C\)
stays inside \(C\): a vertex \(z\) on a path from \(x\in C\) to \(y\in C\)
satisfies \(x\leadsto z\leadsto y\leadsto x\), so \(z\in C\).

For a strong component \(C\) put \(F(C)=\max\{f(x):x\in C\}\).

\emph{Lemma A.}  If \(G\) has an edge \(a\to b\) with \(a\in C\),
\(b\in C'\) and \(C\neq C'\), then \(F(C)>F(C')\).  Let \(x\) be the first
vertex of \(C\cup C'\) to turn grey.  At that moment every other vertex of
\(C\cup C'\) is white.
\emph{Case \(x\in C\).}  Every \(y\in C\) is reached from \(x\) by a path
inside \(C\) (Fact~4), and every \(y\in C'\) by a path running inside \(C\) to
\(a\), across the edge to \(b\), and inside \(C'\) to \(y\).  All vertices of
these paths other than \(x\) are white, so by Fact~3 every
\(y\in(C\cup C')\setminus\{x\}\) has \(f(y)<f(x)\).  Hence
\(F(C)=f(x)>F(C')\).
\emph{Case \(x\in C'\).}  As before, Fact~3 gives \(f(y)\leq f(x)\) for every
\(y\in C'\), so \(F(C')=f(x)\).  No vertex of \(C\) is reachable from \(x\):
if \(x\leadsto z\) with \(z\in C\), then also \(z\leadsto a\to b\leadsto x\),
and \(x\) would lie in \(C\).  By Fact~2 no vertex of \(C\) turns grey during
the interval of \(\mathrm{Visit}(x)\).  They were white at its start, so they
are still white at its end and finish later: \(f(y)>f(x)\) for every
\(y\in C\).  Hence \(F(C)>f(x)=F(C')\).

\emph{Consequence of Lemma A.}  If \(y\leadsto z\) in \(G\) and
\(C(y)\neq C(z)\), then \(F(C(y))>F(C(z))\).  Along the path, an edge inside
one component leaves \(F\) of the current component unchanged, an edge between
distinct components lowers it strictly by Lemma~A, and at least one edge of the
path joins distinct components because the end components differ.

\emph{The second pass.}  We show by induction on the number of trees already
built that, before each new tree is started, the set \(U\) of visited
vertices is a union of strong components, and that the tree started at \(r\)
is exactly \(C(r)\).  Initially \(U=\varnothing\).  Suppose the claim holds
for the earlier trees and a tree is started at an unvisited vertex \(r\).
Because \(U\) is a union of components, all of \(C(r)\) is unvisited.  The
vertices are taken in decreasing order of \(f\), and every vertex taken
earlier is visited (it was either already visited or the root of a tree).
Hence every unvisited vertex has finishing number at most \(f(r)\), and
\(F(C(r))=f(r)\).

The facts for the first pass hold for \(\mathrm{Visit}^{\mathrm R}\) on
\(G^{\mathrm R}\) with ``unvisited'' in place of ``white'', by the same proofs.
So the tree of \(r\) consists of vertices reachable from \(r\) in
\(G^{\mathrm R}\) (Fact~2), and contains every vertex reached from \(r\) in
\(G^{\mathrm R}\) by a path of unvisited vertices (Fact~3).  Each \(y\in C(r)\)
is reached from \(r\) in \(G^{\mathrm R}\) by the reverse of a path from \(y\)
to \(r\) in \(G\), which lies inside \(C(r)\) by Fact~4 and so consists of
unvisited vertices.  Hence \(C(r)\) lies in the tree.  Conversely, let \(y\)
be in the tree.  Then \(y\) is reachable from \(r\) in \(G^{\mathrm R}\), so
\(y\leadsto r\) in \(G\).  If \(C(y)\neq C(r)\), the consequence of Lemma~A
gives \(F(C(y))>F(C(r))=f(r)\), so some \(z\in C(y)\) has \(f(z)>f(r)\).
Such a \(z\) was taken earlier and is visited, so \(C(y)\subseteq U\) because
\(U\) is a union of components, and \(y\) was visited before this tree began.
A call is never made on a visited vertex, so \(y\) is not in the tree, a
contradiction.  Therefore the tree is exactly \(C(r)\), and
\(U\cup C(r)\) is again a union of components.  Every vertex is taken at some
point and is then visited, so every component is a tree.

\emph{Topological numbering.}  The tree with number \(k\) is the component
\(C_k\) of its root \(r_k\), and \(F(C_k)=f(r_k)\) by the paragraph above.
Roots are taken in decreasing order of \(f\), so \(F(C_1)>F(C_2)>\cdots\).
If \(G\) has an edge \(X\to Y\) with \(C(X)\neq C(Y)\), Lemma~A gives
\(F(C(X))>F(C(Y))\), hence \(\operatorname{idx}(X)<\operatorname{idx}(Y)\).

\emph{Cost.}  In each pass every vertex is the argument of exactly one call
(Fact~1 and its analogue), and a call on \(u\) scans the adjacency list of
\(u\) once, doing \(O(1)\) work per entry with the explicit stack.  The
adjacency lists of \(G\) together have \(E\) entries, and so do those of
\(G^{\mathrm R}\).  Colouring, the outer loops and the list \(\mathit{fin}\)
cost \(O(V)\).  The second pass reads \(\mathit{fin}\) backwards, so no sorting
is needed.  The total is \(O(V+E)\).
\end{proof}
\end{annexitem}

\begin{annexitem}{annex:p02-w02-p02-019-2sat}{thm:p02-w02-p02-019-two-sat}{The 2-SAT test}
\begin{proof}
\emph{Clause 1.}  Each edge of \(G_S\) comes from a clause of \(S\).  An edge
\(\overline L\to K\) comes from a clause \(\{L,K\}\), or from the unit clause
\(\{L\}\) with \(K=L\).  Suppose \(v\) satisfies \(S\) and \(v\entails\overline L\).
Then \(L\) is false; the clause \(\{L,K\}\) (or \(\{L\}\)) has a true
literal, which must be \(K\); so \(v\entails K\).  Conversely, suppose every
edge is respected by \(v\), and let \(C\) be a clause.  If \(C=\{L,K\}\) and
\(L\) is false, then \(\overline L\) is true and the edge \(\overline L\to K\)
makes \(K\) true.  If \(C=\{L\}\) and \(L\) is false, then \(\overline L\) is
true and the edge \(\overline L\to L\) makes \(L\) true, which is impossible;
so \(L\) is true.  Every clause is satisfied.  By induction on length, a
valuation that respects every edge also respects every path: if
\(X\leadsto Y\) and \(v\entails X\), then \(v\entails Y\).

\emph{Clause 2.}  The edges contributed by a clause \(\{L,K\}\) are
\(\overline L\to K\) and \(\overline K\to L\).  The companion of the first,
\(\overline K\to\overline{\overline L}\), is the second, since
\(\overline{\overline L}=L\), and conversely.  For a unit clause \(\{L\}\)
the edge \(\overline L\to L\) is its own companion.  So \(X\to Y\) is an edge
iff \(\overline Y\to\overline X\) is.  Reversing a path and complementing its
vertices gives a path, so \(X\leadsto Y\) iff
\(\overline Y\leadsto\overline X\).  Hence \(Z\in C(X)\), that is
\(X\leadsto Z\leadsto X\), holds iff \(\overline X\leadsto\overline Z\leadsto
\overline X\), iff \(\overline Z\in C(\overline X)\).

\emph{Clause 3.}  Let \(C(p_q)=C(\neg p_q)\), so \(p_q\leadsto\neg p_q\) and
\(\neg p_q\leadsto p_q\).  Let \(v\) satisfy \(S\).  By clause~1, \(v\)
respects both paths.  If \(v(p_q)=\mathbf{T}\), the first path makes \(\neg p_q\)
true; if \(v(p_q)=\mathbf{F}\), then \(\neg p_q\) is true and the second path makes
\(p_q\) true.  Both are impossible, so no valuation satisfies \(S\).

\emph{Clause 4.}  Write \(h(X)=\operatorname{idx}(X)\).  By
Theorem~\ref{thm:p02-w02-p02-019-scc}, \(h\) is constant on each component,
different on different components, and \(h(X)\leq h(Y)\) for every edge
\(X\to Y\) (with equality iff the edge lies inside one component).  By
hypothesis \(h(p_q)\neq h(\neg p_q)\) for every \(q\).  For every literal
\(X\) over \(p_1,\ldots,p_n\),
\[
 v_S\entails X\quad\Longleftrightarrow\quad h(X)>h(\overline X).
\]
For \(X=p_q\) this is the definition of \(v_S\).  For \(X=\neg p_q\):
\(v_S\entails\neg p_q\) iff not \(h(p_q)>h(\neg p_q)\), iff
\(h(\neg p_q)>h(p_q)\), because the two values differ.

Let \(X\to Y\) be an edge with \(v_S\entails X\), and suppose
\(v_S\not\entails Y\).  Then \(h(\overline X)<h(X)\) and \(h(Y)<h(\overline
Y)\).  The edge gives \(h(X)\leq h(Y)\), and its companion
\(\overline Y\to\overline X\) (clause~2) gives
\(h(\overline Y)\leq h(\overline X)\).  Together,
\[
 h(\overline X)<h(X)\leq h(Y)<h(\overline Y)\leq h(\overline X),
\]
which is impossible.  So \(v_S\) respects every edge, and by clause~1 it
satisfies \(S\).  This is the argument of Annex~\ref{annex:p04.dh.two-sat}
for the particular rank \(h=\operatorname{idx}\).

\emph{Clause 5.}  Step~(0) is correct, since the empty clause is false under
every valuation.  If step~(3) answers ``unsatisfiable'', some
\(\operatorname{idx}(p_q)=\operatorname{idx}(\neg p_q)\); by
Theorem~\ref{thm:p02-w02-p02-019-scc} the two literals then lie in one
component, and \(S\) is unsatisfiable by clause~3.  If step~(3) passes, no
opposite pair shares a component and \(v_S\) satisfies \(S\) by clause~4.  So
the answer is correct in both cases and a satisfying valuation accompanies
every positive answer.

Cost: step~(0) costs \(O(m)\).  Step~(1) creates \(2n\) empty adjacency lists
for each graph and, for each clause, computes two complements by changing a
binary digit and appends two edges to each graph: \(O(n+m)\).  Step~(2) costs
\(O(V+E)=O(2n+2m)\) by Theorem~\ref{thm:p02-w02-p02-019-scc}.  Steps~(3)
and~(4) compare and assign \(n\) pairs of numbers: \(O(n)\).  The total is
\(O(n+m)\); with no empty clause, \(m\leq N\) and \(n\leq N\), so it is
\(O(N)\).

For the empty list, \(n=m=0\): the graph is empty, step~(3) checks nothing,
and step~(4) returns the valuation making every variable false, whose
restriction to the empty set of occurring variables is the empty assignment;
it satisfies the empty list, as every valuation does.
\end{proof}
\end{annexitem}

\begin{annexitem}{annex:p02-w02-p02-019-xor}{thm:p02-w02-p02-019-xor}{Elimination decides XOR-SAT}
\begin{proof}
Identify a valuation \(v\) with the vector \(x\in\{0,1\}^n\) of
Convention~\ref{def:p02-w02-p02-019-parity-input}; then \(v\) satisfies
equation \(j\) iff \(a_{j1}x_1+\cdots+a_{jn}x_n=b_j\) holds in the
two-element field, because a sum of zeros and ones modulo \(2\) is the parity
of the number of ones.  So satisfying valuations of the system correspond
exactly to its solutions \(x\).

\emph{Row operations preserve solutions.}  Exchanging two rows does not change
the set of equations.  Replacing row \(i'\) by its sum with row \(r\), where
\(i'\neq r\): if \(x\) solves rows \(i'\) and \(r\), adding the two equations
shows that \(x\) solves the new row \(i'\); conversely, if \(x\) solves the
new row \(i'\) and row \(r\), adding them gives back row \(i'\), because
\(y+y=0\) for every entry \(y\).  The other rows are unchanged.  By induction
on the number of operations, the system after the loop has the same solutions
as the input.

\emph{Loop invariant.}  Say that the columns \(1,\ldots,c\) have been
\emph{processed} after the pass for \(c\), and let \(t=r-1\) be the number of
pivots recorded so far, at columns \(c_1<\cdots<c_t\).  The invariant is:
\begin{enumerate}
  \item[(a)] for each \(s\leq t\), row \(s\) has \(a_{sc_s}=1\), and
        \(a_{jc_s}=0\) for every row \(j\neq s\);
  \item[(b)] every row \(j\geq r\) has \(a_{ji}=0\) for every processed
        column \(i\).
\end{enumerate}
Before the loop no column is processed and \(t=0\), so both parts hold
vacuously.  Suppose they hold after the pass for \(c-1\), and consider the
pass for \(c\).  If no row \(i\geq r\) has \(a_{ic}=1\), nothing changes;
(a) still holds, and rows \(j\geq r\) are now zero in column \(c\) as well,
so (b) holds with \(c\) processed.  Otherwise step 2 exchanges two rows with
indices at least \(r\); this leaves rows \(1,\ldots,t\) in place and keeps
(b) for columns before \(c\), so (a) is unaffected and the new row \(r\) has
zeros in columns \(1,\ldots,c-1\) and \(a_{rc}=1\).  In step 3, adding row
\(r\) to another row changes that row only in columns where row \(r\) has a
\(1\), all of which are at least \(c\).  Hence every entry in columns
\(1,\ldots,c-1\) is unchanged, which keeps the pattern of (a) in the earlier
pivot columns \(c_s<c\), and every row other than \(r\) ends with a \(0\) in
column \(c\), since exactly the rows with a \(1\) there received row \(r\).
Thus (a) holds for the \(t+1\) pivots including \((r,c)\).  Rows \(j\geq r+1\)
are zero in columns before \(c\), which were unchanged, and in column \(c\),
so (b) holds after \(r\) is increased by one.  This completes the induction.

\emph{Bound on \(t\).}  Each pivot uses a different row and a different
column, so \(t\leq\min\{m,n\}\).

\emph{Correctness of the answer.}  After the pass for \(c=n\), all columns
are processed, so by (b) every row \(j>t\) has all coefficients zero and
reads \(0=b_j\).  If some such \(b_j=1\), the final system has no solution;
by the preservation of solutions neither has the input, and the answer
``unsatisfiable'' is correct.  Otherwise let \(x^{*}\) be the vector of
\(v^{*}\): \(x^{*}_{c_s}=b_s\) for \(s\leq t\) and \(x^{*}_i=0\) for every
column \(i\) that is not a pivot column.  Consider row \(s\leq t\).  Its
coefficient at \(c_s\) is \(1\); by (a), its coefficient at every other pivot
column \(c_{s'}\) is \(0\); its coefficients at non-pivot columns multiply
zeros.  So its left side evaluates to \(x^{*}_{c_s}=b_s\), and the equation
holds.  Every row \(j>t\) reads \(0=0\).  Hence \(x^{*}\) solves the final
system, and so the input, which proves the claim that \(v^{*}\) satisfies
every equation.  Since the method answers ``satisfiable'' exactly when no row
\(j>t\) has \(b_j=1\), and in that case a satisfying valuation exists, while
in the other case none exists, the answer is correct in both cases.  This
proves item 1.

\emph{Field operations.}  Field additions occur only in step 3.  Each pivot
runs step 3 once, adding row \(r\) to at most \(m-1\) other rows, each by
\(n+1\) entrywise additions.  So there are at most \(t(m-1)(n+1)\) field
operations.  When \(m,n\geq 1\), \(t\leq\min\{m,n\}\) and \(n+1\leq 2n\)
give at most \(2mn\min\{m,n\}\), which is item 2.

\emph{Other steps.}  For each of the \(n\) columns, step 1 inspects at most
\(m\) entries, and step 3 inspects the entry in column \(c\) of each of the
\(m\) rows to select the rows to change: \(O(mn)\) steps in all.  Step 2
exchanges two pointers and step 4 appends one pair, \(O(1)\) steps per
pivot.  Loop control costs \(O(n)\).  The final test inspects at most \(m\)
right-hand sides, and writing \(v^{*}\) costs \(O(n+t)\) steps: initialize
all \(n\) values to false and set one value per pivot.  The total besides
the field operations is \(O(mn+m+n)\), which with item 2 gives item 3.

\emph{Degenerate sizes.}  If \(m=0\), there is no row, every search fails,
\(t=0\), no row \(j>t\) exists, and the answer is ``satisfiable'' with the
all-false valuation; the empty system is indeed satisfied by every valuation.
If \(n=0\), the loop does not run, \(t=0\), and every row reads \(0=b_j\);
the answer is ``unsatisfiable'' iff some \(b_j=1\).  This is item 4.
\end{proof}
\end{annexitem}

\begin{annexitem}{annex:p02-w02-p02-019-tt}{thm:p02-w02-p02-019-truth-table}{Truth tables decide}
\begin{proof}
\emph{Reduction to finitely many assignments.}  Let \(X=\{p_1,\ldots,p_k\}\)
be the jointly occurring variables.  Every valuation \(v\) agrees on \(X\)
with \(u^{\circ}\) for \(u=v\mathbin{\upharpoonright}X\), and both agree on
every variable occurring in each \(\varphi_i\).  By
Lemma~\ref{lem:prop-coincidence}, \(\widehat v(\varphi_i)=
\widehat{u^{\circ}}(\varphi_i)\) for every \(i\).  Conversely each \(u\) is the
restriction of the valuation \(u^{\circ}\).  So the vectors of values
\((\widehat v(\varphi_1),\ldots,\widehat v(\varphi_r))\) taken over all
valuations \(v\) are exactly those taken over the \(2^k\) assignments \(u\).

\emph{Correctness.}  By this, some valuation makes \(\varphi_1\) true iff some
\(u\) does; every valuation does iff every \(u\) does; every valuation gives
\(\varphi_1,\varphi_2\) the same value iff every \(u\) does, which is the
definition of \(\varphi_1\equiv\varphi_2\) in
\cref{sec:prop-syntax-semantics-02}; and a countervaluation to
\(\{\varphi_1,\ldots,\varphi_{r-1}\}\entails\varphi_r\), a valuation
satisfying the premises and falsifying the conclusion, exists iff some \(u\)
gives \(\mathbf{T}\) to the premises and \(\mathbf{F}\) to \(\varphi_r\).  When
\(r=1\) the premise set is empty, and the answer coincides with validity, as
\(\varnothing\entails\varphi\) iff \(\entails\varphi\).

\emph{The evaluation scan.}  Write \(\neg^j\chi\) for \(\chi\) preceded by
\(j\) negation signs.  We show by induction on the formation of \(\chi\): if
the scan begins to read \(\chi\) when the operator stack is \(O\) followed by
\(j\geq 0\) copies of \(\neg\), where \(O\) is empty or ends in a symbol other
than \(\neg\), and the value stack is \(W\), then after reading \(\chi\) the
operator stack is \(O\) and the value stack is \(W\) followed by the single
value \(\widehat{u^{\circ}}(\neg^j\chi)\).
For \(\chi=p_q\) or \(\chi=\bottom\), the scan pushes \(u^{\circ}(p_q)\) or
\(\mathbf{F}\), the value the truth clauses give, and absorbs the \(j\) waiting
negations, each by the negation clause; absorption stops at the last symbol
of \(O\), which is not \(\neg\).
For \(\chi=\neg\chi'\), the scan pushes \(\neg\), so it begins to read
\(\chi'\) with \(j+1\) waiting negations; by the induction hypothesis it ends
with operator stack \(O\) and the value of \(\neg^{j+1}\chi'=\neg^j\chi\).
For \(\chi=(\chi_1\circ\chi_2)\), the scan pushes the opening parenthesis,
which is not \(\neg\); by the induction hypothesis with no waiting negation,
reading \(\chi_1\) adds its value \(b_1\).  It then pushes \(\circ\), and
reading \(\chi_2\) adds \(b_2\) in the same way.  At the closing parenthesis it
pops \(b_2\), \(\circ\), \(b_1\) and the parenthesis, pushes the value of
\(b_1\circ b_2\) given by the truth clause for \(\circ\), which is
\(\widehat{u^{\circ}}(\chi)\), and absorbs the \(j\) negations now on top.
Applied to a whole formula with empty stacks and \(j=0\), the scan leaves
exactly its value.  Each symbol is pushed at most once and popped at most
once, so one scan of \(\chi\) costs \(O(|\chi|)\).

\emph{Bound.}  The binary counter steps through the \(2^k\) assignments,
changing \(O(k)\) bits per step in the worst case; for \(k\geq 1\),
\(k\leq\ell\) because every variable of \(X\) occurs.  For each assignment the
scans cost \(O(\ell)\) in total, and the comparison of the \(r\) values
\(O(r)\subseteq O(\ell)\).  So the run costs \(O(2^k\ell)\); for
\(k=0\) there is one assignment and the cost is \(O(\ell)\).  Since
\(2^k\ell\) is finite for every input, the procedure terminates.  For a single
formula, \(\ell=|\varphi|\).
\end{proof}
\end{annexitem}

\begin{annexitem}{annex:p02-w02-p02-019-np}{thm:p02-w02-p02-019-sat-in-np}{SAT is in NP under the fixed encoding}
\begin{proof}
Let \(x\) have length \(n\).  A machine deciding \(R(x,y)\) proceeds as follows.

\emph{1.  Decode and recognize.}  Cut \(x\) into fixed-width blocks and
translate each into its character, rejecting if the length is not a multiple
of the block width or some block codes no character: \(O(n)\) steps.  Check
that the resulting character string is a formula: replace each variable
token, the letter \texttt{p} followed by binary digits and the terminator
\texttt{;}, and each occurrence of \(\bottom\), by a marker \(\mathtt A\); then
repeat a pass that replaces every \(\neg\mathtt A\) by \(\mathtt A\) and
every \((\mathtt A\circ\mathtt A)\) with \(\circ\in\{\land,\lor,\to\}\) by
\(\mathtt A\), copying to a second tape, until a pass changes nothing.  Call
a string a \emph{marked formula} if it arises from a formula by replacing
some subformula occurrences by \(\mathtt A\).  In a marked formula, a
\(\neg\) begins the subformula it governs, and a subformula that begins with
\(\mathtt A\) is that marker; so an occurrence of \(\neg\mathtt A\) or of
\((\mathtt A\circ\mathtt A)\) is a subformula occurrence, and replacing it
gives a marked formula.  Conversely, if the result of such a replacement is a
marked formula, so was the string before it, by the formation clauses.  A
marked formula other than \(\mathtt A\) contains such a pattern: take a
subformula occurrence of least length that is not a marker; its immediate
subformulas are shorter, so they are markers, all atomic ones having been
replaced at the start.  Hence the string
is a formula iff the passes end with the single marker \(\mathtt A\).  Each
pass that does not stop
shortens the string, so there are at most \(n\) passes of \(O(n)\) steps:
\(O(n^2)\).  This polynomial bound is all that is used; the linear-time
recognition noted in Definition~\ref{def:p02-sat-taut} is not needed.

\emph{2.  Number the variables.}  For each variable token, compare its digit
string with those of all earlier tokens to decide whether it is the first
occurrence of its variable, and count the earlier first occurrences to obtain
its rank \(i\), so that the token names \(p_{a_i}\).  There are at most \(n\)
tokens, each compared with at most \(n\) others, each comparison taking
\(O(n)\) symbols with \(O(n)\) head movement per symbol: \(O(n^4)\) steps.  The
number \(k\) of first occurrences is computed on the way; reject unless
\(|y|=k\).

\emph{3.  Evaluate.}  Rewrite the formula with each token of rank \(i\)
replaced by the \(i\)-th bit of \(y\) and \(\bottom\) by \(0\), then apply
passes that replace every \(\neg b\) by \(1-b\) and every \((b\circ c)\) by the
value of the truth table of \(\circ\) at \(b,c\).  As in step~1 there are at
most \(n\) passes of \(O(n)\) steps, ending in one bit, which is
\(\widehat{v_y}(\varphi)\) by the truth clauses.  Accept iff it is \(1\).

The total is \(O(n^4)\), a polynomial in \(n\), and \(R\) is the predicate of
the statement.  Every accepted \(y\) has length \(k\); distinct variables
have disjoint first-occurrence tokens of at least one block each, so
\(k\leq n\).

Now \(x\in\mathsf{SAT}\) iff \(x\) encodes a formula \(\varphi\) and some
valuation \(v\) satisfies \(\varphi\).  Given such \(v\), let \(y\) have
\(i\)-th bit \(1\) iff \(v(p_{a_i})=\mathbf{T}\).  Then \(v_y\) agrees with \(v\) on
every variable of \(\varphi\), so \(v_y\entails\varphi\) by
Lemma~\ref{lem:prop-coincidence}, \(R(x,y)\) holds and \(|y|=k\leq n\).
Conversely, if \(R(x,y)\) holds, then \(x\) encodes \(\varphi\) and \(v_y\)
satisfies it, so \(x\in\mathsf{SAT}\).  A malformed \(x\) is rejected in
step~1 for every \(y\), and it is not in \(\mathsf{SAT}\) by
Definition~\ref{def:p02-sat-taut}.  With \(q(n)=n\) and this \(R\), the
condition of Convention~\ref{def:p02-complexity-standing} holds, so
\(\mathsf{SAT}\in\mathsf{NP}\).
\end{proof}
\end{annexitem}

\begin{annexitem}{annex:p02-w02-p02-019-red}{thm:p02-w02-p02-019-reductions}{Complements and composites of reductions}
\begin{proof}
\emph{Clause 1.}  For every \(x\), \(x\in A\) iff \(f(x)\in B\).  Negating
both sides, \(x\in\overline A\) iff \(f(x)\in\overline B\).  The map and its
running time are unchanged.

\emph{Clause 2.}  Let \(f\) and \(g\) be computed within \(t_f(n)\) and
\(t_g(n)\) steps, polynomials which we may take nondecreasing.  A machine
writes at most one output symbol per step, so \(|f(x)|\leq t_f(|x|)\).
Computing \(f(x)\) and then \(g(f(x))\) takes at most
\(t_f(n)+t_g(t_f(n))\) steps plus the \(O(t_f(n))\) steps needed to move the
intermediate string to an input position, a polynomial in \(n\).  And
\(x\in A\) iff \(f(x)\in B\) iff \(g(f(x))\in D\).

\emph{Clause 3.}  Let \(f\) reduce \(A\) to \(B\) within \(t_f\), and let
\(B\) have polynomial \(q\) and polynomial-time predicate \(R_B\) as in
Convention~\ref{def:p02-complexity-standing}, with \(q\) nondecreasing.  Put
\(R_A(x,y)\) iff \(R_B(f(x),y)\), and \(q_A(n)=q(t_f(n))\).  Then \(R_A\) is
decided in polynomial time by computing \(f(x)\) and running the machine for
\(R_B\), whose input has length at most \(t_f(n)+|y|\).  If \(x\in A\), then
\(f(x)\in B\), so some \(y\) with \(|y|\leq q(|f(x)|)\leq q_A(n)\) satisfies
\(R_B(f(x),y)\).  If some \(y\) satisfies \(R_A(x,y)\), then \(f(x)\in B\)
and \(x\in A\).  So \(A\in\mathsf{NP}\).
\end{proof}
\end{annexitem}

\begin{annexitem}{annex:p02-w02-p02-019-neg}{thm:p02-w02-p02-019-negation-maps}{Negation maps}
\begin{proof}
By the negation clause, \(v\entails\neg\varphi\) iff \(v\not\entails\varphi\).
Hence every valuation satisfies \(\varphi\) iff no valuation satisfies
\(\neg\varphi\): \(\varphi\) is valid iff \(\neg\varphi\) is unsatisfiable.

By the formation clauses, \(\neg\varphi\) is the string \(\neg\) followed by
\(\varphi\), so \(\nu\) prefixes one fixed block.  If \(x\) encodes
\(\varphi\): \(x\in\mathsf{TAUT}\) iff \(\varphi\) is valid iff \(\neg\varphi\)
is unsatisfiable iff \(\nu(x)\notin\mathsf{SAT}\), since \(\nu(x)\) is a
well-formed encoding.  So \(x\in\mathsf{TAUT}\) iff
\(g_1(x)\in\overline{\mathsf{SAT}}\), and \(x\in\overline{\mathsf{SAT}}\)
iff \(\varphi\) is unsatisfiable iff \(\neg\varphi\) is valid iff
\(g_2(x)\in\mathsf{TAUT}\).  If \(x\) is malformed: \(x\notin\mathsf{TAUT}\),
and \(g_1(x)\), the encoding of \(p_0\), is in \(\mathsf{SAT}\) because the
valuation making \(p_0\) true satisfies it, so
\(g_1(x)\notin\overline{\mathsf{SAT}}\); and \(x\in\overline{\mathsf{SAT}}\),
while \(g_2(x)\), the encoding of \(\neg\bottom\), is in \(\mathsf{TAUT}\)
because \(\widehat v(\bottom)=\mathbf{F}\) and so \(\widehat v(\neg\bottom)=\mathbf{T}\)
for every \(v\).  Both maps are therefore many-one reductions as stated.
Each decides well-formedness in \(O(n^2)\) steps as in step~1 of
Annex~\ref{annex:p02-w02-p02-019-np} and then writes \(x\) with one
block in front, or a fixed string: polynomial time.
\end{proof}
\end{annexitem}

\begin{annexitem}{annex:p02-w02-p02-019-taut}{thm:p02-w02-p02-019-taut-conp}{TAUT is coNP-complete}
\begin{proof}
\emph{Membership.}  By Lemma~\ref{thm:p02-w02-p02-019-negation-maps}, \(g_1\)
reduces \(\mathsf{TAUT}\) to \(\overline{\mathsf{SAT}}\), so by
Lemma~\ref{thm:p02-w02-p02-019-reductions}, clause~1, it reduces
\(\overline{\mathsf{TAUT}}\) to \(\mathsf{SAT}\).  Since
\(\mathsf{SAT}\in\mathsf{NP}\) (Theorem~\ref{thm:p02-w02-p02-019-sat-in-np}),
clause~3 of the same lemma gives \(\overline{\mathsf{TAUT}}\in\mathsf{NP}\).
So \(\mathsf{TAUT}\) is the complement of an \(\mathsf{NP}\) problem, that is,
\(\mathsf{TAUT}\in\mathsf{coNP}\) by Convention~\ref{def:p02-complexity-standing}.

\emph{Hardness.}  Let \(L\in\mathsf{coNP}\).  Then \(L=\overline K\) for some
\(K\in\mathsf{NP}\), and \(\overline L=K\).  By the hardness half of
Theorem~\ref{thm:p02-sat-np-complete} there is a polynomial-time reduction
\(f\) of \(K\) to \(\mathsf{SAT}\).  By
Lemma~\ref{thm:p02-w02-p02-019-reductions}, clause~1, \(f\) reduces
\(\overline K=L\) to \(\overline{\mathsf{SAT}}\), and \(g_2\) reduces
\(\overline{\mathsf{SAT}}\) to \(\mathsf{TAUT}\)
(Lemma~\ref{thm:p02-w02-p02-019-negation-maps}).  By clause~2 of the
reduction lemma, \(g_2\circ f\) is a polynomial-time reduction of \(L\) to
\(\mathsf{TAUT}\).  Membership and hardness together are
\(\mathsf{coNP}\)-completeness under polynomial-time many-one reductions.
\end{proof}
\end{annexitem}

\begin{annexitem}{annex:p02-w02-p02-019-npco}{thm:p02-w02-p02-019-np-conp}{Validity, equivalence and NP}
\begin{proof}
Let \(H\) be \(\mathsf{TAUT}\) or the equivalence language of
Theorem~\ref{thm:p14-dc-finite-cpl}.  Both are \(\mathsf{coNP}\)-complete:
the first by Theorem~\ref{thm:p02-w02-p02-019-taut-conp}, the second by
Theorem~\ref{thm:p14-dc-finite-cpl}, whose hardness proof reduces validity to
equivalence.

Suppose \(H\in\mathsf{NP}\).  Every \(L\in\mathsf{coNP}\) reduces to \(H\), so
\(L\in\mathsf{NP}\) by Lemma~\ref{thm:p02-w02-p02-019-reductions}, clause~3;
thus \(\mathsf{coNP}\subseteq\mathsf{NP}\).  If \(K\in\mathsf{NP}\), then
\(\overline K\in\mathsf{coNP}\subseteq\mathsf{NP}\), so
\(K=\overline{\overline K}\) is the complement of an \(\mathsf{NP}\) problem
and \(K\in\mathsf{coNP}\); thus \(\mathsf{NP}\subseteq\mathsf{coNP}\).  Hence
\(\mathsf{NP}=\mathsf{coNP}\).  Conversely, if \(\mathsf{NP}=\mathsf{coNP}\),
then \(H\in\mathsf{coNP}=\mathsf{NP}\).
\end{proof}
\end{annexitem}
